\documentclass[preprint,11pt]{elsarticle}
\usepackage[T1]{fontenc}\usepackage[utf8]{inputenc}\usepackage{lmodern}
\usepackage{amsmath,amssymb,amsthm,booktabs,graphicx,tabularx,multirow,array,longtable,xcolor}
\usepackage{microtype}\emergencystretch=3em
\usepackage[hidelinks]{hyperref}
\usepackage{cleveref}
\usepackage{xr-hyper}
\externaldocument{main}
\newtheorem{proposition}{Proposition}
% generated by code/experiments/aggregate.py -- do not edit
\newcommand{\ATwoK}{7}
\newcommand{\ATwoN}{7}
\newcommand{\BudgetLsOne}{27}
\newcommand{\BudgetLsTen}{36}
\newcommand{\BudgetLsTenth}{19}
\newcommand{\BudgetMainLsTenth}{9}
\newcommand{\BudgetMainNlTenth}{4}
\newcommand{\BudgetNlOne}{28}
\newcommand{\BudgetNlTen}{36}
\newcommand{\BudgetNlTenth}{10}
\newcommand{\ClsLsAtFivetwelve}{12}
\newcommand{\ClsLsAtOnetwentyeight}{7}
\newcommand{\ClsLsAtSixtyfour}{8}
\newcommand{\ClsLsAtTentwentyfour}{13}
\newcommand{\ClsLsAtThirtytwo}{5}
\newcommand{\ClsLsAtTwofiftysix}{9}
\newcommand{\ClsNlAtFivetwelve}{5}
\newcommand{\ClsNlAtOnetwentyeight}{8}
\newcommand{\ClsNlAtSixtyfour}{5}
\newcommand{\ClsNlAtTentwentyfour}{8}
\newcommand{\ClsNlAtThirtytwo}{5}
\newcommand{\ClsNlAtTwofiftysix}{4}
\newcommand{\ClsTieAtFivetwelve}{7}
\newcommand{\ClsTieAtOnetwentyeight}{9}
\newcommand{\ClsTieAtSixtyfour}{11}
\newcommand{\ClsTieAtTentwentyfour}{3}
\newcommand{\ClsTieAtThirtytwo}{14}
\newcommand{\ClsTieAtTwofiftysix}{11}
\newcommand{\CpuDirectHi}{139}
\newcommand{\CpuDirectLo}{116}
\newcommand{\CpuGp}{44}
\newcommand{\CpuPodmlp}{86}
\newcommand{\DrawAgree}{83\%}
\newcommand{\DrawSign}{99\%}
\newcommand{\ExpBChanged}{15}
\newcommand{\ExpBImproved}{90\%}
\newcommand{\ExpBRetention}{88\%}
\newcommand{\ExpBShift}{0.24}
\newcommand{\ExpBTotal}{48}
\newcommand{\ExpFGain}{0.47}
\newcommand{\ExpFKept}{5}
\newcommand{\ExpFN}{7}
\newcommand{\ExtMedianExt}{-0.03}
\newcommand{\ExtMedianInterp}{0.16}
\newcommand{\FomMedianMs}{6.0}
\newcommand{\FomN}{36}
\newcommand{\FullCellsN}{28}
\newcommand{\FullOverPod}{0.95}
\newcommand{\HfourAB}{10}
\newcommand{\HfourAC}{1}
\newcommand{\HfourAP}{0.006}
\newcommand{\HoneAMedian}{-0.059}
\newcommand{\HoneAP}{0.874}
\newcommand{\HoneBK}{5}
\newcommand{\HtwoPHolm}{0.145}
\newcommand{\IsoOverArd}{1.92}
\newcommand{\KrrNarrowN}{28}
\newcommand{\KrrNarrowRatio}{13}
\newcommand{\KrrNarrowShareThree}{18\%}
\newcommand{\KrrWideEightBeat}{61\%}
\newcommand{\KrrWideEightMed}{6.5}
\newcommand{\KrrWideEightShare}{21\%}
\newcommand{\KrrWideThirtytwoBeat}{64\%}
\newcommand{\KrrWideThirtytwoMed}{2.7}
\newcommand{\KrrWideThirtytwoShare}{57\%}
\newcommand{\LeakErrorRatio}{229}
\newcommand{\LeakMaxMatches}{0}
\newcommand{\LeakPlantedFlagged}{20}
\newcommand{\MetricAgreeMed}{95\%}
\newcommand{\NMainProblems}{24}
\newcommand{\NOwnProblems}{36}
\newcommand{\NlClearTieBandZeroonefive}{6}
\newcommand{\NlClearTieBandZerothree}{3}
\newcommand{\NlClearTieBandZerozerofive}{10}
\newcommand{\NnFloorF}{$2.0\times10^{-4}$}
\newcommand{\PLeFourLs}{20}
\newcommand{\PLeFourN}{30}
\newcommand{\PLeFourNl}{8}
\newcommand{\PoneP}{0.18}
\newcommand{\PoneRho}{-0.38}
\newcommand{\PooledRmseBase}{1.07}
\newcommand{\PooledRmseCv}{0.62}
\newcommand{\PtwoGain}{0.60}
\newcommand{\PtwoPHolm}{0.051}
\newcommand{\SeedAllLs}{80\%}
\newcommand{\SeedAllNl}{57\%}
\newcommand{\SoneCorrect}{6}
\newcommand{\SoneN}{10}

\providecommand{\tightlist}{\setlength{\itemsep}{0pt}\setlength{\parskip}{0pt}}
\DeclareUnicodeCharacter{2212}{\ensuremath{-}}\DeclareUnicodeCharacter{0233}{\ensuremath{\bar y}}\DeclareUnicodeCharacter{2265}{\ensuremath{\ge}}\DeclareUnicodeCharacter{2192}{\ensuremath{\to}}\DeclareUnicodeCharacter{03C1}{\ensuremath{\rho}}\DeclareUnicodeCharacter{2264}{\ensuremath{\le}}\DeclareUnicodeCharacter{00B7}{\ensuremath{\cdot}}
\usepackage{calc}
\newcommand{\real}[1]{#1}
\renewcommand{\thesection}{S\arabic{section}}\renewcommand{\thetable}{S\arabic{table}}\renewcommand{\thefigure}{S\arabic{figure}}\renewcommand{\theequation}{S\arabic{equation}}
\begin{document}
\begin{frontmatter}
\title{Supplementary Material: Choosing a Fast Stand-In for a Slow Simulation}
\author{[Author names to be supplied by the authors]}
\end{frontmatter}

\section{Proofs}\label{sec:S-proofs}
Both results are standard; we include short proofs for completeness.
\paragraph{Proof of Proposition~\ref{prop:loo}.} Let $f^{-i}$ minimise $\sum_{j\ne i}(y_j-f_j)^2+\lambda N\lVert f\rVert_K^2$ and let $y'$ equal $y$ except $y'_i=f^{-i}_i$. For any $f$, the objective with data $y'$ equals the leave-one-out objective plus $(y'_i-f_i)^2\ge0$, so $f^{-i}$, which makes the extra term zero and minimises the first, also minimises the objective with data $y'$. By uniqueness of the ridge solution, $f^{-i}=Hy'$. Hence $f^{-i}_i=(Hy')_i=\hat y_i-H_{ii}y_i+H_{ii}f^{-i}_i$, so $(1-H_{ii})f^{-i}_i=\hat y_i-H_{ii}y_i$ and $y_i-f^{-i}_i=(y_i-\hat y_i)/(1-H_{ii})$. For $K=Q\,\mathrm{diag}(w)Q^\top$ with $Q$ orthogonal, $H=Q\,\mathrm{diag}\big(w/(w+\lambda N)\big)Q^\top$ and, since $\sum_jQ_{ij}^2=1$, $1-H_{ii}=\sum_jQ_{ij}^2\,\lambda N/(w_j+\lambda N)$. For the polynomial model with a thin design matrix $F=U\,\mathrm{diag}(s)V^\top$, $1-H_{ii}=(1-\sum_jU_{ij}^2)+\sum_jU_{ij}^2\lambda N/(s_j^2+\lambda N)$. A bug in an early version of the tool (Sec.~\ref{sec:S-deviations}, V9) came from computing $1-H_{ii}$ by subtraction and clipping it. \qed
\paragraph{Proof of Proposition~\ref{prop:floor}.} Let $P=V_rV_r^\top$. For $x\in\bar{\mathbf u}+\mathrm{span}(V_r)$, $x-\mathbf u=(x-\bar{\mathbf u}-P(\mathbf u-\bar{\mathbf u}))-(I-P)(\mathbf u-\bar{\mathbf u})$, where the first term lies in $\mathrm{span}(V_r)$ and the second in its orthogonal complement. By Pythagoras, $\lVert x-\mathbf u\rVert_2^2\ge\lVert(I-P)(\mathbf u-\bar{\mathbf u})\rVert_2^2$. Dividing by $\lVert\mathbf u\rVert_2$ and averaging over runs gives the claim. \qed

\section{Pre-specification record}\label{sec:S-prereg}
Hypotheses, success criteria and analysis code were committed to the project's git repository before the results they concern (a self-hosted history, not an independent registry). Table~\ref{tab:S-prereg} lists, for each experiment, the protocol commit and the first commit that contains its results; the check is automatic (\texttt{audit/prereg\_order\_check.py}) and shows that no protocol was committed after its results. It shows protocol-before-results in history, not protocol-before-first-run. Experiments H, G3 and the Stage-6 experiments T, C and S are labelled post hoc in motivation: they were designed in response to earlier results or to reviews (Sec.~\ref{sec:S-deviations}).
\begin{table}[h]\centering\caption{Order of pre-specification and results commits.}\label{tab:S-prereg}\small\resizebox{\textwidth}{!}{\begin{tabular}{lllrl}
\toprule
Exp & Protocol commit & First results commit & Hours before & Order ok \\
\midrule
A & bcf7f30 & 2d0e2d4 & 9.33 & yes \\
D & bcf7f30 & 2abe263 & 10.99 & yes \\
B & bcf7f30 & 72cb02c & 11.07 & yes \\
A2 & 72cb02c & 703ad59 & 2.38 & yes \\
E & 72cb02c & 703ad59 & 2.38 & yes \\
A3 & c9a7b93 & 39b1b7f & 3.66 & yes \\
G & c9a7b93 & 39b1b7f & 3.66 & yes \\
F & c9a7b93 & 39b1b7f & 3.66 & yes \\
H8 & 48a0698 & 84f2a21 & 0.08 & yes \\
H32 & 48a0698 & 84f2a21 & 0.08 & yes \\
G3 & 4cd411a & b37dd4e & 0.05 & yes \\
6T & a99c2c6 & 0e45aae & 32.58 & yes \\
6C & a99c2c6 & 1dc799e & 53.43 & yes \\
6S & a99c2c6 & 9ebaf04 & 56.97 & yes \\
R & 2619b07 & a88ad5f & 19.76 & yes \\
7T & a096980 & af431a1 & 0.51 & yes \\
7C & a096980 & 7c4953c & 1.51 & yes \\
7S & a096980 & 36abd58 & 2.01 & yes \\
7N & a096980 & 6fc0b43 & 2.86 & yes \\
7SP & a096980 & b22bb2e & 3.46 & yes \\
7E & a096980 & 5e1f2d3 & 5.85 & yes \\
\bottomrule
\end{tabular}
}\end{table}

\section{Earlier-stage hypotheses and verdicts}\label{sec:S-earlier}
Table~\ref{tab:S-earlier} lists the hypotheses and primary tests of the earlier stages (original 16 problems; independent replication on eight new problems; mechanism controls; longer-training and precision checks). Most confirmatory tests of the first stage were not supported; the one supported at cell level (unequal tuning) rests on 5 of 16 problems.
\begin{table}[h]\centering\caption{Earlier-stage tests. $\bar y$ as in \eqref{eq:y}.}\label{tab:S-earlier}\footnotesize\resizebox{\textwidth}{!}{\begin{tabular}{p{1.5cm}p{4.6cm}p{3.3cm}p{1.7cm}p{3.0cm}}
\toprule
ID & Criterion & Observed & $p$ & Verdict \\
\midrule
H1a (v1) & median $\bar y$ at $N{=}32$ positive (reduced basis better) & median -0.059 & 0.874 & not supported \\
H1b (v1) & $\bar y$ falls with $N$ in most problems & 5/16 problems & 0.962 & not supported \\
H2 (v1) & pilot features beat $\{\log_2N,p\}$ & CI [-0.07, 6.40] & 0.145 (Holm) & not supported \\
H4a (v1) & unequal tuning moves wins toward the more tuned side & 10 vs 1 discordant cells & 0.006 & supported (cell level) \\
P1 (Exp A) & Spearman($f_1,\bar y$)$<0$, 8 new problems & $\rho=-0.38$ & 0.180 & not rejected \\
P2 (Exp A) & frozen rule beats baseline out of sample & gain 0.60 & 0.051 (Holm) & not rejected (knife-edge) \\
Q1/Q2 (A2) & $\ge4$ of 7 smooth candidates with $\bar y\ge1$; none network-clear & 7/7; 0 & -- & met (confounded with $p$) \\
D1 & narrow-range kernel ridge within $3\times$ of GP in $\ge75\%$ of cells & median ratio 14.3; 13\% within $3\times$ & -- & rejected \\
E & iso-GP vs ARD-GP and kernel ridge & ratios 1.92, 8.1 & -- & marginal-likelihood rule met (low $p$ only) \\
G1 & full-field GP within $3\times$ of POD+GP & median ratio 0.95 & -- & supported \\
H (post hoc) & wide-range kernel ridge within $3\times$ ($B{=}8$ or 32) & medians 6.5, 2.7 & -- & neither rule met (mixed) \\
B1/B2 & 4$\times$ epochs: sign kept in $\ge85\%$; median shift $<0.5$ & 88\%; 0.24 & -- & met, but 15 of 48 classes change \\
F & float64 + L-BFGS: $\ge5$ of 7 decisive gaps remain $\ge1$ decade & 5 of 7 & -- & met \\
\bottomrule
\end{tabular}
}\end{table}
\paragraph{Pilot diagnostic} A frozen rule built from the original problems was tested on eight new problems from four new families: $\rho=\PoneRho$ (exact $p=\PoneP$); out-of-sample gain \PtwoGain{} with Holm-adjusted $p=\PtwoPHolm$ (10$^6$ bootstrap draws); the threshold rule classified \SoneCorrect{} of \SoneN{} non-tie cells correctly. The pilot statistic also correlates with the number of inputs, and all problems with $f_1<-2$ among the 24 have three inputs; we conclude that the diagnostic is inconclusive and do not use it in the tool.
\paragraph{Robustness} Class labels agree between halves of the random-search draws in \DrawAgree{} of cells (sign agreement when both clear: \DrawSign); \SeedAllLs{} of reduced-basis-clear cells and \SeedAllNl{} of network-clear cells have all three seeds beyond the tie band; class labels computed from the median instead of the mean error agree in \MetricAgreeMed{} of cells; the extrapolation test set gives a median $\bar y$ of \ExtMedianExt{} against \ExtMedianInterp{} on the interpolation set at $N=1024$.

\section{Deviations from the plan}\label{sec:S-deviations}
Every change to the plan is recorded with its date and reason in the repository (\texttt{research\_design/deviations.md} for the first stage and \texttt{deviations\_v2.md} afterwards); the text is reproduced here.
\hypertarget{deviations-disclosures-and-post-hoc-exploratory-analyses}{%
\subsection*{Deviations, disclosures and post-hoc (exploratory) analyses}\label{deviations-disclosures-and-post-hoc-exploratory-analyses}}

\begin{longtable}[]{@{}
  >{\raggedright\arraybackslash}p{(\columnwidth - 8\tabcolsep) * \real{0.2000}}
  >{\raggedright\arraybackslash}p{(\columnwidth - 8\tabcolsep) * \real{0.2000}}
  >{\raggedright\arraybackslash}p{(\columnwidth - 8\tabcolsep) * \real{0.2000}}
  >{\raggedright\arraybackslash}p{(\columnwidth - 8\tabcolsep) * \real{0.2000}}
  >{\raggedright\arraybackslash}p{(\columnwidth - 8\tabcolsep) * \real{0.2000}}@{}}
\toprule\noalign{}
\begin{minipage}[b]{\linewidth}\raggedright
\#
\end{minipage} & \begin{minipage}[b]{\linewidth}\raggedright
Item
\end{minipage} & \begin{minipage}[b]{\linewidth}\raggedright
Type
\end{minipage} & \begin{minipage}[b]{\linewidth}\raggedright
Detail
\end{minipage} & \begin{minipage}[b]{\linewidth}\raggedright
Effect on confirmatory results
\end{minipage} \\
\midrule\noalign{}
\endhead
\bottomrule\noalign{}
\endlastfoot
D1 & Run restarts & operational & First launch killed (BLAS thread oversubscription) before any job finished; the container restarted twice later and killed in-flight jobs. Completed jobs are atomic (JSON written only on completion); the runner resumes and skips them. Identical code/protocol. & none \\
D2 & Interim analysis & disclosure & A premature analysis was run when 47 of 48 jobs existed (seeds 0,1 complete; my file count wrongly included \texttt{environment.json}). It used only complete seed batches (0,1) and was \textbf{not acted upon}; no code, hypothesis or setting changed. Output kept in \texttt{results/\allowbreak{}interim\_\allowbreak{}2seed\_\allowbreak{}disclosed/\allowbreak{}}. The confirmatory analysis is the 3-seed run. Verdicts of H1a, H1b, H2, H4a were identical in interim and final. & none \\
D3 & Seeds & protocol & Seeds 0--2 (mandatory) completed; extension seeds 3--4 \textbf{not run} (compute/time). & declared in protocol \\
D4 & Post-hoc exploratory analyses & exploratory & Added after seeing confirmatory results; every one is labelled \emph{exploratory} in the manuscript and none changes a confirmatory verdict. See \texttt{results/\allowbreak{}exploratory/\allowbreak{}}. & none \\
D5 & Non-reproducible search sampling (bug) & disclosure + fix & Found in the fresh-clone audit: \texttt{hpo.py} seeded the random-search RNG with Python's per-process-salted \texttt{hash(method)}, so trials 1-7 of every hyper-parameter search are sampled differently in every process. The experiment ran with this code. Effect: a from-scratch re-execution samples different trial configurations than those recorded (statistically equivalent, not bit-identical). All sampled configs and their validation/test errors are stored in \texttt{results/\allowbreak{}raw/\allowbreak{}*.json}; \texttt{code/\allowbreak{}experiments/\allowbreak{}replay\_\allowbreak{}check.py} refits stored best configs and reproduces 20/20 recorded test errors exactly (relative diff 0.0). Fixed afterwards with \texttt{zlib.crc32} (regression test added). Results were \textbf{not} regenerated. & none on reported numbers; re-runs will differ numerically \\
\end{longtable}

\hypertarget{deviations-disclosures-and-corrections-in-the-v2-v2.1-v2.2-cycle}{%
\subsection*{Deviations, disclosures and corrections in the v2 / v2.1 / v2.2 cycle}\label{deviations-disclosures-and-corrections-in-the-v2-v2.1-v2.2-cycle}}

\begin{longtable}[]{@{}
  >{\raggedright\arraybackslash}p{(\columnwidth - 8\tabcolsep) * \real{0.2000}}
  >{\raggedright\arraybackslash}p{(\columnwidth - 8\tabcolsep) * \real{0.2000}}
  >{\raggedright\arraybackslash}p{(\columnwidth - 8\tabcolsep) * \real{0.2000}}
  >{\raggedright\arraybackslash}p{(\columnwidth - 8\tabcolsep) * \real{0.2000}}
  >{\raggedright\arraybackslash}p{(\columnwidth - 8\tabcolsep) * \real{0.2000}}@{}}
\toprule\noalign{}
\begin{minipage}[b]{\linewidth}\raggedright
\#
\end{minipage} & \begin{minipage}[b]{\linewidth}\raggedright
Item
\end{minipage} & \begin{minipage}[b]{\linewidth}\raggedright
Type
\end{minipage} & \begin{minipage}[b]{\linewidth}\raggedright
Detail
\end{minipage} & \begin{minipage}[b]{\linewidth}\raggedright
Effect on pre-registered verdicts
\end{minipage} \\
\midrule\noalign{}
\endhead
\bottomrule\noalign{}
\endlastfoot
V1 & Exp B shift sign convention & bug found by independent audit (F1), fixed \textbf{before Exp B was analysed} & \texttt{shift\ =\ y6000\ −\ y1500} is negative when NL improves; B2 ``shift toward NL \textless{} 0.5'' would have passed trivially. Replaced by \texttt{shift\_\allowbreak{}toward\_\allowbreak{}NL\ =\ y1500\ −\ y6000}; unit test added. & none (B not yet analysed) \\
V2 & A2 baseline p0 miscounted & factual error in addendum v2.1 found by independent audit, verified & p0 stated 2/20; correct 3/20 (burgers\_L1 f1 = −1.9972, ȳ = 2.25). Amended in addendum v2.2 \textbf{before A2 ran}: p0 = 0.15, success k ≥ 4 of 7. & none (A2 not yet run) \\
V3 & Diagnostic verdict wording & re-characterisation & ``Not replicated'' → \emph{pre-registered criteria not met; inconclusive; carried by one family (nlheat 88.3 \% of P2 gain; ρ = +0.49 without it); confounded with p (Spearman 0.708)}. & none (numbers unchanged) \\
V4 & Stage chronology & disclosure & Stage 3 (addendum v2.1) was committed within minutes after the Exp A and Exp D results; Stage 4 (addendum v2.2) after two independent audits. All preregistrations are on a self-hosted git history; git tags could not be pushed to the remote. & none \\
V5 & Shared search draws & design property found by audit (F2/F3) & The random-search RNG depends only on (seed, method): the same 7 sampled configs are used for every instance and N within a seed (3 distinct sets per method in the whole study). v1 additionally used a process-salted hash (D5). & none on verdicts; changes how ``tuning'' must be described; sensitivity analysis planned \\
V6 & Neural batch/step count & design property found by audit (F02) & batch = N for N ≤ 512 (1500 steps) but 512 for N = 1024 (3000 steps): the neural step count doubles at N = 1024. Exp F removes this (full batch). & none on verdicts \\
V7 & Test/doc fixes & housekeeping & stale test count fixed (36 instances); weaker leakage tests replaced; BLAS thread env set at the top of entry scripts. & none \\
\end{longtable}

\hypertarget{v9-surroselect-smoke-test-bug-before-any-registered-stage-6-run-2026-11-02}{%
\subsubsection*{V9 --- SurroSelect smoke-test bug (before any registered Stage-6 run) --- 2026-11-02}\label{v9-surroselect-smoke-test-bug-before-any-registered-stage-6-run-2026-11-02}}

During a smoke test of the tool on the new problem \texttt{lshape\_\allowbreak{}L1} (seed 0, N = 64 and 256) and a quick check of the KRR and polynomial members on \texttt{heatcomp\_\allowbreak{}L2} and \texttt{convdiff\_\allowbreak{}L1} (N = 64, 256), the wide-range kernel-ridge member returned a test error of 0.5 (10x worse than any other member). Cause: the closed-form leave-one-out residual divides by 1 - H\_ii, which was clipped at 1e-6 (and computed by cancellation), producing spuriously tiny LOO scores for near-interpolating fits; the length-scale grid was also defined for inputs on {[}0,1{]} but with the {[}-1,1{]} limits of the benchmark's wide-range control. Fix (committed before any registered run): 1 - H\_ii computed without cancellation (KRR: sum\_j Q\_ij\^{}2 lam N/(w\_j + lam N); polynomial: (1 - sum U\_ij\^{}2) + sum U\_ij\^{}2 lam N/(s\_j\^{}2 + lam N)), no clipping above 1e-12, ridge values with condition number above 1e12 skipped, length-scale grid {[}0.05, 15{]}·sqrt(p) on {[}0,1{]}-scaled inputs (= {[}0.1, 30{]}·sqrt(p) on {[}-1,1{]}); a regression test was added. The smoke tests showed test errors of all pool members on lshape\_L1 (seed 0, N = 64, 256) and of the KRR and polynomial members on three problems; the pool, criteria, tie band, folds and all pre-registered thresholds were NOT changed after seeing them. Addendum v3.0 stated that no surrogate error on the new problems had been seen; this entry corrects that statement.

\hypertarget{v10-surroselect-pickling-bug-found-by-the-plain-file-validation-2026-11-04}{%
\subsubsection*{V10 --- SurroSelect pickling bug found by the plain-file validation (2026-11-04)}\label{v10-surroselect-pickling-bug-found-by-the-plain-file-validation-2026-11-04}}

After the \texttt{metric} option was added (v3.1), a fitted selector stored a lambda and could no longer be saved by \texttt{surroselect\ fit\ -\/\allowbreak{}-model}; the command-line unit test had been run before the change and the full tool test suite was not re-run. The plain-file validation (experiment E8) exposed the failure (``Can't get local object''). Fix: the metric function is a property computed from the metric name. The registered evaluation (experiments C, S, R) uses the library interface without pickling and is unaffected; all tool tests were re-run after the fix.


\section{Data-leakage audit}\label{sec:S-leak}
Training, test and extrapolation inputs are drawn from separate named random streams. Over all \NAllProblems{} problems the number of exact input matches between the training pool and the test or extrapolation sets is \LeakMaxMatches. A planted-leak test (20 test samples copied into the training set of problem \texttt{heatcomp\_L1}, $N=64$) is flagged by the overlap check (\LeakPlantedFlagged{} exact matches, none without the leak), and the error on the leaked samples is \LeakErrorRatio{} times lower than on the other test samples, showing that a leak would be visible. Validation errors do not depend on test labels (unit test, corrupting all test labels). Fitting steps (POD, scalers, hyper-parameters, cross-validation folds) use training runs only. Two properties remain: training subsets are nested within a seed, and seeds share the pool and test set, so seed variance understates sampling variance at large $N$. Comparisons in the mechanism study that use the best network by \emph{test} error are oracle comparisons and are labelled as such.

\section{Per-problem results}\label{sec:S-perproblem}
Table~\ref{tab:S-perproblem} gives $\bar y$ (reduced basis against direct networks, validation-selected members, seed means) for every problem; positive values favour the reduced basis.
{\footnotesize\begin{longtable}{lrrrr}
\caption{$\bar y$ per problem and budget (positive: reduced basis better).}\label{tab:S-perproblem}\\
\toprule
Problem & $p$ & $N{=}64$ & $N{=}256$ & $N{=}1024$ \\
\midrule
\endhead
\texttt{advdiff\_L1} & 4 & -0.04 & 0.17 & 0.21 \\
\texttt{advdiff\_L2} & 4 & -0.09 & -0.00 & -0.07 \\
\texttt{bratu\_L1} & 3 & 2.78 & 2.98 & 3.31 \\
\texttt{bratu\_L2} & 3 & 2.14 & 2.97 & 3.08 \\
\texttt{burgers\_L1} & 4 & -0.06 & 0.68 & 2.25 \\
\texttt{burgers\_L2} & 4 & 0.01 & 0.08 & 0.11 \\
\texttt{convdiff\_L1} & 3 & -0.03 & 0.60 & 0.48 \\
\texttt{convdiff\_L2} & 4 & -0.05 & 0.10 & 0.04 \\
\texttt{darcy\_L1} & 4 & 0.15 & 1.12 & 1.13 \\
\texttt{darcy\_L2} & 10 & -0.00 & -0.05 & -0.19 \\
\texttt{darcyhd\_L1} & 16 & 0.04 & -0.02 & -0.13 \\
\texttt{darcyhd\_L2} & 32 & 0.03 & 0.00 & -0.07 \\
\texttt{darcylin\_L1} & 16 & 0.18 & 0.63 & 0.48 \\
\texttt{darcylin\_L2} & 32 & 0.04 & 0.43 & 0.18 \\
\texttt{diff1d\_L1} & 4 & 0.89 & 2.36 & 2.70 \\
\texttt{diff1d\_L2} & 8 & 0.17 & 0.62 & 1.00 \\
\texttt{fisher\_L1} & 4 & -0.21 & 0.09 & 0.52 \\
\texttt{fisher\_L2} & 4 & -0.26 & -0.26 & -0.44 \\
\texttt{heat2d\_L1} & 4 & 0.18 & 0.93 & 1.38 \\
\texttt{heat2d\_L2} & 4 & -0.04 & -0.10 & -0.34 \\
\texttt{heatcomp\_L1} & 4 & 0.00 & 0.52 & 0.88 \\
\texttt{heatcomp\_L2} & 9 & -0.01 & -0.06 & -0.12 \\
\texttt{heatmove\_L1} & 3 & 0.13 & 0.13 & -0.17 \\
\texttt{heatmove\_L2} & 5 & -0.04 & -0.02 & -0.15 \\
\texttt{heatmv3\_L1} & 3 & 0.12 & -0.11 & -0.44 \\
\texttt{heatp\_L1} & 3 & 0.77 & 2.49 & 2.87 \\
\texttt{heatp\_L2} & 3 & 1.08 & 2.71 & 2.90 \\
\texttt{helm\_L1} & 3 & 1.20 & 2.69 & 2.93 \\
\texttt{helm\_L2} & 3 & -0.16 & 2.08 & 2.64 \\
\texttt{incl3\_L1} & 3 & 0.02 & 0.05 & -0.21 \\
\texttt{incl\_L1} & 4 & -0.02 & -0.00 & -0.25 \\
\texttt{incl\_L2} & 4 & 0.02 & -0.01 & -0.13 \\
\texttt{kdv\_L1} & 4 & -0.14 & -0.26 & -0.23 \\
\texttt{kdv\_L2} & 4 & -0.04 & -0.03 & 0.29 \\
\texttt{lshape\_L1} & 4 & 0.20 & 1.14 & 1.66 \\
\texttt{lshape\_L2} & 7 & -0.11 & 0.01 & -0.01 \\
\texttt{nlheat\_L1} & 3 & 1.25 & 2.51 & 2.72 \\
\texttt{nlheat\_L2} & 3 & 0.75 & 2.16 & 2.58 \\
\texttt{reacdiff\_L1} & 3 & 0.85 & 1.62 & 2.57 \\
\texttt{reacdiff\_L2} & 3 & 0.57 & 1.96 & 2.38 \\
\texttt{transport\_L1} & 4 & -0.26 & 0.05 & 0.75 \\
\texttt{transport\_L2} & 4 & -0.32 & -0.23 & -0.33 \\
\texttt{wave\_L1} & 4 & 0.16 & 0.29 & 0.78 \\
\texttt{wave\_L2} & 4 & 0.10 & -0.06 & -0.04 \\
\bottomrule
\end{longtable}
}

\section{Descriptive run on the earlier problems}\label{sec:S-retro}
Addendum v3.0 (Sec.~6.4) promised a descriptive run of the tool on the earlier problems, with one seed and $N=256$. The run was not reported in the first version of this work and was made afterwards (this is deviation V11 in the log). It is \emph{design-informed}: the pool and the grids of the tool were chosen with knowledge of these problems, so the run is not evidence for the criteria. On the \RetroN{} earlier problems solved with our own code, the median regret against the oracle (best pool member by test error) was \RetroMedian{} (largest \RetroMax; \RetroShare{} of the problems within a factor two); the largest regrets were on \RetroWorst. The tool's error was a median \RetroLeg{} times that of the benchmark's validation-selected method (better by more than 10\% on \RetroBetter{} problems, worse on \RetroWorse), and \RetroGp{} times that of the benchmark's POD+GP. The basis flag was raised on \RetroFlagged{} of the problems.

\section{The tool}\label{sec:S-tool}
\paragraph{Installation and use.} The package \texttt{surroselect} (directory \texttt{surroselect/} of the repository; MIT licence, copyright holder to be set by the authors) depends on NumPy, SciPy and scikit-learn; PyTorch and Optuna are needed only for the optional \texttt{TunedMLP}. In Python:
\begin{verbatim}
sel = SurrogateSelector(cv=4).fit(X, Y)    # X: runs x inputs, Y: runs x outputs
sel.learning_curve()
print(make_report(sel, target=0.01))
Y_new = sel.predict(X_new)
\end{verbatim}
On the command line, with plain \texttt{.npz} files that hold arrays \texttt{X} and \texttt{Y}:
\begin{verbatim}
surroselect fit train.npz --model m.pkl --target 0.01 --report report.md
surroselect predict m.pkl X_new.npy --out Y_pred.npy
surroselect evaluate m.pkl test.npz
\end{verbatim}
\paragraph{Default pool.} POD+GP (Mat\'ern-5/2, rank 32), POD+GP (squared exponential, rank 32), POD+GP (Mat\'ern-5/2, rank 8), POD+KRR with a wide grid of kernel, length-scale and ridge selected by closed-form leave-one-out error, POD+polynomial (degree up to six, at most 3000 features, leave-one-out), and a GP on the full field (Mat\'ern-5/2 or squared exponential; only for $m\le2000$). An optional metric, the normalised RMSE $\lVert Y-\hat Y\rVert_F/\lVert Y-\bar Y_{\rm train}\rVert_F$, is available for outputs with a large offset.
\paragraph{Validation of the tool on plain files.} To check the tool itself, without the benchmark pipeline in the loop, three of the new problems (\texttt{heatcomp\_L1}, \texttt{lshape\_L1}, \texttt{convdiff\_L1}; $N=256$) were exported to \texttt{.npz} files and processed with the three command-line commands as separate processes. All \PlainOk{} of \PlainN{} sequences ran without error, and the error printed by \texttt{evaluate} agreed with an independent computation from the saved predictions to within $10^{-12}$ in \PlainAgree{} of \PlainN{} cases (\texttt{results/s6/plain\_files\_validation.json}).
\paragraph{Tests.} \NToolTests{} unit tests for the tool (closed-form leave-one-out against brute force and its stability near interpolation, exact polynomial recovery, selection accuracy, the basis-limit flag, the power-law inversion, the report, the metric option and the command-line round trip), \NTpTests{} tests for the finite-element problems and the \NBenchTests{} tests of the benchmark package are collected by the repository; \texttt{results/logs/tests\_*.log} records the final runs. A fitted selector could not be saved in an early version after the metric option was added (deviation V10); the plain-file validation found it.
\paragraph{Known limits.} The tool assumes exchangeable runs (A3); it does not model noise; for $p\gtrsim10$ the default GP kernels have many hyper-parameters fitted on at most 300 runs; the forecast is an extrapolation without guarantee; and the optional neural member uses a small search budget.

\end{document}
